%========================================================
% SECTION 6.4 — CHSH Nonlocality
%========================================================

\newpage

\section{CHSH Nonlocality}

%--------------------------------------------------------
% Subsection 6.4.1 — Objective and Physical Motivation
%--------------------------------------------------------

\subsection*{Objective and Physical Motivation}

This experiment investigates the nonlocal correlations of the bipartite quantum state through the Clauser--Horne--Shimony--Holt (CHSH) parameter.

The analysis compares two physically distinct propagation mechanisms:

\begin{itemize}
    \item deterministic phase evolution generated by local unitary transformations;
    \item depolarizing noise described by an effective two-qubit quantum channel.
\end{itemize}

The objective is to determine how these mechanisms affect the observable violation of the CHSH inequality and to distinguish between:

\begin{itemize}
    \item apparent reductions of CHSH violation caused by measurement-axis misalignment;
    \item genuine loss of nonlocality caused by decoherence.
\end{itemize}

To this end, two different CHSH quantities are evaluated:

\begin{itemize}
    \item \(S_{\mathrm{fixed}}\), computed using fixed laboratory measurement axes;
    \item \(S_{\max}\), computed from the correlation tensor and corresponding to the maximal CHSH value achievable after optimization over measurement settings.
\end{itemize}

This distinction is physically important because a decrease of \(S_{\mathrm{fixed}}\) does not necessarily imply degradation of the intrinsic nonlocal properties of the quantum state.

%--------------------------------------------------------
% Subsection 6.4.2 — Model and Analytical Predictions
%--------------------------------------------------------

\subsection*{Model and Analytical Predictions}

The reference state is the Bell state

\[
|\Psi^+\rangle
=
\frac{1}{\sqrt{2}}
\left(
|01\rangle + |10\rangle
\right).
\]

Two different state families are analyzed.

\medskip

\textbf{Phase-evolved states.}

The first family is generated through deterministic local phase evolution:

\begin{equation}
|\psi(\theta)\rangle
=
\frac{1}{\sqrt{2}}
\left(
|01\rangle + e^{i\theta}|10\rangle
\right),
\label{eq:chsh_phase_state}
\end{equation}

obtained by applying the local unitary transformation

\[
U_A(\theta)
\otimes
I_B.
\]

\medskip

\textbf{Depolarized states.}

The second family corresponds to the depolarizing channel:

\begin{equation}
\rho(p)
=
(1-p)\rho_0
+
\frac{p}{4}I_4,
\label{eq:chsh_depolarized_state}
\end{equation}

where

\[
\rho_0
=
|\Psi^+\rangle
\langle\Psi^+|,
\qquad
p \in [0,1].
\]

\medskip

The fixed CHSH measurement settings are chosen to be optimal for the Bell state at \( \theta = 0 \), whose correlation tensor is

\[
T(0)
=
\begin{pmatrix}
1 & 0 & 0 \\
0 & 1 & 0 \\
0 & 0 & -1
\end{pmatrix}.
\]

The local analyzer directions for subsystem \(A\) are

\[
\mathbf{a}_1
=
\begin{pmatrix}
1 \\ 0 \\ 0
\end{pmatrix},
\qquad
\mathbf{a}_2
=
\begin{pmatrix}
0 \\ 1 \\ 0
\end{pmatrix},
\]

while the analyzer directions for subsystem \(B\) are

\[
\mathbf{b}_1
=
\frac{1}{\sqrt{2}}
\begin{pmatrix}
1 \\ 1 \\ 0
\end{pmatrix},
\qquad
\mathbf{b}_2
=
\frac{1}{\sqrt{2}}
\begin{pmatrix}
1 \\ -1 \\ 0
\end{pmatrix}.
\]

\begin{figure}[htbp]
    \centering
    \includegraphics[width=\linewidth]{figures/experiment_4/experiment_4_bloch_axes.png}
    \caption{
    Fixed CHSH measurement axes represented on the Bloch spheres of subsystems \(A\) and \(B\).
    The analyzer directions define the fixed laboratory configuration used to evaluate \(S_{\mathrm{fixed}}\).
    }
    \label{fig:exp4_bloch_axes}
\end{figure}

The corresponding Bloch-sphere geometry is shown in Figure~\ref{fig:exp4_bloch_axes}.

\medskip

With these settings, \(S_{\mathrm{fixed}}\) represents the CHSH value measured using a fixed analyzer configuration, whereas \(S_{\max}\) represents the maximal violation achievable after optimization over local measurement axes.

\medskip

For the phase-evolved Bell state, the correlation tensor becomes

\begin{equation}
T(\theta)
=
\begin{pmatrix}
\cos\theta & -\sin\theta & 0 \\
\sin\theta & \cos\theta & 0 \\
0 & 0 & -1
\end{pmatrix}.
\label{eq:chsh_phase_tensor}
\end{equation}

Using the fixed analyzer directions introduced above, the corresponding CHSH value is

\[
S_{\mathrm{fixed}}(\theta)
=
2\sqrt{2}\cos\theta.
\]

In the numerical plots, the quantity

\[
|S_{\mathrm{fixed}}|
\]

is represented, since CHSH violation depends on the condition

\[
|S| > 2.
\]

The maximal CHSH value is instead independent of the phase parameter:

\[
S_{\max}
=
2\sqrt{2}.
\]

This result reflects the invariance of nonlocality under local unitary transformations.

\medskip

For the depolarizing channel, the correlation tensor is uniformly contracted:

\[
T(p)
=
(1-p)T(0).
\]

Consequently,

\[
S_{\mathrm{fixed}}(p)
=
2\sqrt{2}(1-p),
\qquad
S_{\max}(p)
=
2\sqrt{2}(1-p).
\]

The CHSH violation threshold is determined from

\[
S_{\max}(p)
>
2,
\]

which gives

\[
p
<
p_{\mathrm{CHSH}}
=
1-\frac{1}{\sqrt{2}}
\simeq
0.2929.
\]

%--------------------------------------------------------
% Subsection 6.4.3 — Numerical Results
%--------------------------------------------------------

\subsection*{Numerical Results}

The numerical results are summarized in Figure~\ref{fig:exp4_chsh}.

\begin{figure}[htbp]
    \centering
    \includegraphics[width=\linewidth]{figures/experiment_4/experiment_4_chsh_nonlocality_summary.png}
    \caption{
    CHSH behavior under deterministic phase evolution and depolarizing noise.
    The left panel shows \( |S_{\mathrm{fixed}}| \) and \( S_{\max} \) as functions of the phase parameter \( \theta \).
    The right panel shows the corresponding quantities as functions of the depolarizing parameter \( p \).
    The red dotted line indicates the classical CHSH bound \( |S|=2 \), while the purple dotted line corresponds to the Tsirelson bound \(2\sqrt{2}\).
    }
    \label{fig:exp4_chsh}
\end{figure}

In the deterministic phase-evolution regime, \( |S_{\mathrm{fixed}}| \) oscillates between zero and the Tsirelson bound.

This behavior originates from the relative rotation between the fixed laboratory analyzer axes and the phase-dependent correlation tensor.

At the same time, the maximal CHSH value remains constant:

\[
S_{\max}
=
2\sqrt{2},
\]

demonstrating that the intrinsic nonlocality of the state is fully preserved.

\medskip

In the depolarizing regime, the numerical and analytical curves overlap within numerical precision.

Both \( |S_{\mathrm{fixed}}| \) and \( S_{\max} \) decrease linearly with the depolarization parameter \(p\), crossing the classical CHSH bound at

\[
p_{\mathrm{CHSH}}
\simeq
0.2929.
\]

Unlike the deterministic phase case, the reduction of CHSH violation now reflects a genuine loss of nonlocal correlation strength.

%--------------------------------------------------------
% Subsection 6.4.4 — Discussion
%--------------------------------------------------------

\subsection*{Discussion}

The two propagation mechanisms considered in this experiment produce qualitatively different effects on CHSH nonlocality.

In the deterministic phase-evolution regime, the quantum state remains pure and maximally entangled for all values of the phase parameter \( \theta \).

Consequently, the oscillatory behavior of \( |S_{\mathrm{fixed}}| \) does not indicate degradation of nonlocality. Instead, it reflects the fact that the fixed laboratory analyzer directions become progressively misaligned with the rotated correlation tensor.

This explains why \( |S_{\mathrm{fixed}}| \) may fall below the classical CHSH threshold, or even vanish completely, while the optimized quantity \( S_{\max} \) remains fixed at the Tsirelson bound.

The apparent loss of CHSH violation is therefore caused entirely by measurement-axis mismatch.

\medskip

The depolarizing regime exhibits fundamentally different behavior.

The channel reduces the magnitude of the correlation tensor itself, producing a genuine contraction of the available nonlocal correlations.

As a consequence, no optimization of the analyzer directions can recover CHSH violation once the depolarization parameter exceeds the threshold

\[
p_{\mathrm{CHSH}}
=
1-\frac{1}{\sqrt{2}}.
\]

The experiment therefore establishes an important operational distinction:

\begin{itemize}
    \item \(S_{\mathrm{fixed}}\) is sensitive to both measurement misalignment and physical degradation;
    \item \(S_{\max}\) isolates the maximal nonlocality intrinsically present in the quantum state.
\end{itemize}

%--------------------------------------------------------
% Subsection 6.4.5 — Physical Interpretation
%--------------------------------------------------------

\subsection*{Physical Interpretation}

This experiment separates coherent fiber-induced transformations from incoherent decoherence mechanisms.

A deterministic relative phase acts as a local unitary transformation generated by birefringence-induced phase accumulation between polarization modes.

Such a mechanism rotates the correlation tensor but does not modify its singular values. Consequently, the intrinsic nonlocal resources of the Bell state remain fully preserved.

However, a fixed analyzer configuration may fail to detect the CHSH violation if the laboratory axes are no longer aligned with the rotated tensor.

\medskip

Depolarization has a fundamentally different physical interpretation.

It models the effective loss of polarization coherence caused by unresolved degrees of freedom, stochastic fluctuations, frequency-dependent propagation effects, or polarization-mode dispersion.

In this case, the correlation tensor is not merely rotated but genuinely contracted, leading to irreversible degradation of the nonlocal correlations.

The comparison between \(S_{\mathrm{fixed}}\) and \(S_{\max}\) therefore provides an operational diagnostic tool:

\begin{itemize}
    \item if \(S_{\mathrm{fixed}}\) changes while \(S_{\max}\) remains constant, the effect is compatible with coherent unitary evolution;
    \item if \(S_{\max}\) decreases, the quantum state itself has lost nonlocal correlation strength.
\end{itemize}

For the quantum fiber simulator, this distinction is particularly important.

A fiber-induced phase shift may hide CHSH violation from fixed laboratory measurements without destroying the underlying entanglement resource, whereas depolarizing effects correspond to genuine physical degradation of the entangled photon pair.